%-*-latex-*-
%-*-latex-*-
\newcommand\COURSE{ciss350}
\newcommand\ASSESSMENT{q10902}
\newcommand\ASSESSMENTTYPE{Quiz}
\newcommand\POINTS{\textwhite{xxx/xxx}}

\input{myquizpreamble}
\input{yliow}
\input{\COURSE}
%-*-latex-*-

%-*-latex-*-

%-*-latex-*-
\renewcommand\TITLE{\ASSESSMENTTYPE \ \ASSESSMENT}

\renewcommand\EMAIL{}
\renewcommand\AUTHOR{eagraham1@cougars.ccis.edu}

\textwidth=6in
\begin{document}
\topmatterthree

Here's how you describe (in text) for a binary tree with a heap tree structure.
For this tree
\begin{console}[fontsize=\footnotesize]
             (10)
           /      \
        (6)        (7)
       /   \      /   \
    (5)     (2)(0)     (3)
   /
(1)
\end{console}
you write
\begin{console}[fontsize=\footnotesize]
{10, 6, 7, 5, 2, 0, 3, 1}
\end{console}
i.e., list the values using breadth-first traversal (like what we did in class).

%------------------------------------------------------------------------------
\nextq
Execute \verb!build_heap! on the array
$\{7, 5, 3, 1, 2, 4, 6, 8\}$ and describe the resulting maxheap
as an array.
(Note: \verb!build_heap! builds the maxheap \textit{bottom-up}.)
\\
\ANSWER
\begin{answercode}
{8, 7, 6, 5, 2, 4, 3, 1}
\end{answercode}

%------------------------------------------------------------------------------
\nextq
The following is already a maxheap:
\begin{console}[fontsize=\footnotesize]
              (9)
           /       \
        (8)         (7)
       /   \       /   \
    (6)     (3) (1)     (4)
   /   \
(5)     (2)
\end{console}
During heapsort, what is the resulting array after
\textit{four} extract-root operations.
Recall that for heapsort the roots extracted are \textit{not} thrown away!
\\
\ANSWER
\begin{answercode}
{5, 3, 4, 2, 1, 6, 7, 8, 9}
\end{answercode}

%------------------------------------------------------------------------------
\nextq
Suppose your \verb!std::vector< int >! is sorted (in ascending order).
What is the runtime of the \textit{fastest} algorithm
\begin{console}
void ascending_to_maxheap(std::vector< int > & v);
\end{console}
that converts it a maxheap?
\\
\ANSWER
\begin{answerlong}
$T(n) = O(n)$
\end{answerlong}
(You should only have
$O(1)$ or $O(\lg n)$ or $O(n)$ or $O(n \lg n)$ or $O(n^2)$ or ..., i.e.,
it should be of the form $O(n^a \lg^b n)$.)
%------------------------------------------------------------------------------

%------------------------------------------------------------------------------
\nextq
Write a function that checks if a \verb!std::vector< int >! object is a maxheap,
returning the \textit{first} index
(scanning \textit{from index 0} to the end of the array)
where the maxheap property is violated.
If the \verb!std::vector< int >! object is a maxheap, \verb!-1! is returned.
\begin{console}[fontsize=\footnotesize]
int maxheap_violation_index(const std::vector< int > & v);
\end{console}
\\
\ANSWER
\begin{answercode}
int LEFT(int i)
{
  return 2 * i + 1;
}

int RIGHT(int i)
{
  return 2 * i + 2;
}

int maxheap_violation_index(const std::vector< int > & v)
{
  for(int i = 0; i < v.size(); i++)
  {
    int l = LEFT(i);
    int r = RIGHT(i);

    if (l < v.size() && v[i] < v[l])
      return i;

    if (r < v.size() && v[i] < v[r])
      return i;
  }
  return -1;
}
\end{answercode}
%------------------------------------------------------------------------------

%------------------------------------------------------------------------------
%-*-latex-*-
\newpage
\input{instructions}
\end{document}

